\documentclass[../../main.tex]{subfiles}% 注意这里的文件路径不能用 ./main.tex ,否则用latexmk编译子文件会报错
\graphicspath{{\subfix{./image/}}{\subfix{../../image/}}} % 指定图片引用路径,后续可以直接使用图片文件名
% 如果图片存放在上级目录,则使用 ../../image/ 作为备用路径

% 例如:
% \begin{figure}[H]
% \centering
% \includegraphics[width=0.3\textwidth]{图.png}
% \caption{}
% \label{figure:图}
% \end{figure}
% 注意:上述\label{}一定要放在\caption{}之后,否则引用图片序号会只会显示??.

\begin{document}
% SOURCE p=16; printed=1; object=ch00-sec01-title; region=middle
\section{集合论的语言}\label{section:ch00-set-language}
% SOURCE p=16; printed=1; object=ch00-sec01-intro; region=middle
假定读者熟悉集合论的基本概念；以下讨论主要是为了确定我们的术语。

% SOURCE p=16; printed=1; object=ch00-number-systems; region=lower-middle
\textbf{数系。} 我们对基本数系采用如下记号：
\begin{align*}
\mathbb{N}&=\text{正整数集（不包括零）},\\
\mathbb{Z}&=\text{整数集},\\
\mathbb{Q}&=\text{有理数集},\\
\mathbb{R}&=\text{实数集},\\
\mathbb{C}&=\text{复数集}.
\end{align*}
% SOURCE p=16; printed=1; object=ch00-logic-notation; region=bottom
\textbf{逻辑。} 我们将避免使用数理逻辑中的特殊符号，而使表述尽可能接近通常的英语。不过，我们将用缩写 iff 表示“当且仅当”。

% SOURCE p=16,17; printed=1,2; object=ch00-contrapositive; region=bottom,top

\begin{remark}
初等逻辑中有一点，学生常常认识得不够充分：若 $A$ 和 $B$ 是数学命题，而 $-A$、$-B$ 分别是它们的否定，则陈述“$A$ 蕴含 $B$”在逻辑上等价于其逆否命题“$-B$ 蕴含 $-A$”。因此，可以通过假设 $-B$ 并推出 $-A$ 来证明 $A$ 蕴含 $B$，我们将经常这样做。这与反证法不同；反证法是假设 $A$ 和 $-B$ 同时成立，再推出矛盾。
\end{remark}
% SUPPLEMENT: id=v14-ch00-contrapositive-proof; parent=ch00-contrapositive; kind=proof; origin=independent-derivation
\begin{proof}
\textbf{逆否等价的补写证明（重构新增）。}在本书通常采用的经典逻辑下，先假设 $A$ 蕴含 $B$。若 $-B$ 成立而 $A$ 也成立，则由该蕴含关系得到 $B$，与 $-B$ 矛盾，所以 $-A$ 成立。这证明 $-B$ 蕴含 $-A$。反之，假设 $-B$ 蕴含 $-A$，并设 $A$ 成立。若 $-B$ 成立，就会得到 $-A$，与 $A$ 矛盾；因此 $-B$ 不成立，由经典逻辑得到 $B$。这证明 $A$ 蕴含 $B$。逆否法所选择的待证命题是 $-B\Rightarrow-A$；反证法则直接把 $A$ 与 $-B$ 一起作为临时假设并导出矛盾，两者在论证组织方式上不同。
\end{proof}
% END SUPPLEMENT: id=v14-ch00-contrapositive-proof



% SOURCE p=17; printed=2; object=ch00-sets-power-set; region=upper
\textbf{集合。} “族”（family、collection）与“集合”（set）同义，使用这些词通常是为了避免“集合的集合”这样的说法。
\begin{definition}[集合及其子集族]\label{definition:ch00-sets-power-set}
空集记为 $\varnothing$，集合 $X$ 的所有子集构成的族记为 $\mathscr{P}(X)$：
\begin{align*}
\mathscr{P}(X)&=\{E:E\subseteq X\}.
\end{align*}
\end{definition}
% SUPPLEMENT: id=v14-ch00-power-set-remark; parent=ch00-sets-power-set; kind=motivation+hypotheses; origin=independent-derivation
\begin{remark}
\textbf{动机与条件说明（重构新增）。}幂集把“研究 $X$ 中哪些元素”转化为“研究 $X$ 的哪些子集”，因而能把集合族本身作为对象。这里 $X$ 必须是一个集合；$\mathscr{P}(X)$ 的成员是 $X$ 的子集，不能把 $E\subseteq X$ 与 $E\in X$ 混为一谈。采用非严格包含，才能同时把 $X$ 自身和空集纳入幂集：$X\subseteq X$；而空集没有违反成员条件的元素，所以 $\varnothing\subseteq X$。因此 $\mathscr{P}(\varnothing)=\{\varnothing\}$，它并非空集。“空集”和幂集的记号属于约定；这里没有有限性、可数性或额外的代数结构假设。
\end{remark}
% END SUPPLEMENT: id=v14-ch00-power-set-remark


这里以及其他地方，包含符号 $\subseteq$ 均按非严格意义理解；即陈述“$E\subseteq X$”包括 $E=X$ 的可能性。

% SOURCE p=17; printed=2; object=ch00-union-intersection; region=upper-middle

\Needspace{9\baselineskip}
\begin{definition}[集合族的并与交]\label{definition:ch00-union-intersection}
若 $\mathscr{E}$ 是一个集合族，可以构成其中各成员的并与交：
\begin{align*}
\bigcup_{E\in\mathscr{E}}E&=\{x:x\in E\text{ 对某个 }E\in\mathscr{E}\text{ 成立}\},\\
\bigcap_{E\in\mathscr{E}}E&=\{x:x\in E\text{ 对所有 }E\in\mathscr{E}\text{ 成立}\}.
\end{align*}
\end{definition}
% SUPPLEMENT: id=v14-ch00-union-intersection-remark; parent=ch00-union-intersection; kind=motivation+hypotheses; origin=independent-derivation
\begin{remark}
\textbf{动机与条件说明（重构新增）。}并对应“至少满足一个成员条件”，交对应“同时满足每个成员条件”；存在量词与全称量词的区别就是两种运算的区别。集合族的成员必须是集合，才能谈论 $x\in E$。若集合族非空，其交包含在任取的一个成员中，不需要额外指定全集。若集合族为空，并是空集；空族的交则必须在预先固定的全集 $X$ 内解释，此时约定为空族的交等于 $X$，因为每个 $x\in X$ 都满足空的全称条件。没有固定全集时，不能把空族的交解释为“所有集合组成的集合”。以下涉及空指标族的交或补集时，均在已明确的固定全集内讨论。
\end{remark}
% END SUPPLEMENT: id=v14-ch00-union-intersection-remark


% SOURCE p=17; printed=2; object=ch00-indexed-family; region=middle
通常，考虑带指标的集合族更为方便：
\begin{align*}
\mathscr{E}&=\{E_\alpha:\alpha\in A\}=\{E_\alpha\}_{\alpha\in A},
\end{align*}
此时，并与交分别记为
\begin{align*}
\bigcup_{\alpha\in A}E_\alpha,\qquad \bigcap_{\alpha\in A}E_\alpha.
\end{align*}
% SOURCE p=17; printed=2; object=ch00-disjoint-family; region=middle-lower

\begin{definition}[互不相交的集合族]\label{definition:ch00-disjoint-family}
若只要 $\alpha\neq\beta$ 就有 $E_\alpha\cap E_\beta=\varnothing$，则称这些集合 $E_\alpha$ \textbf{互不相交}。“互不相交的集合族”与“由互不相交的集合组成的族”这两种说法可以互换，“集合的不交并”与“互不相交的集合的并”也一样。
\end{definition}
% SUPPLEMENT: id=v14-ch00-disjoint-remark; parent=ch00-disjoint-family; kind=motivation+hypotheses; origin=independent-derivation
\begin{remark}
\textbf{动机与条件说明（重构新增）。}互不相交保证一个元素至多属于一个成员，因此在分块讨论或计数时不会重复。必须限定为不同指标 $\alpha\ne\beta$：若也要求 $E_\alpha\cap E_\alpha=\varnothing$，就会迫使每个成员均为空。这里要求的是每一对不同成员的交为空，强于“所有成员的共同交为空”；例如 $\{1,2\}$、$\{2,3\}$、$\{1,3\}$ 的共同交为空，但任意两者都有公共元素。按指标理解时，同一个非空集合不能在两个不同指标处重复出现；空集可以重复，因为它与任何集合的交仍为空。
\end{remark}
% END SUPPLEMENT: id=v14-ch00-disjoint-remark



% SOURCE p=17; printed=2; object=ch00-set-sequences; region=lower
考虑以 $\mathbb{N}$ 为指标集的集合族时，我们通常采用记号
\begin{align*}
\{E_n\}_{n=1}^{\infty}\quad\text{或}\quad\{E_n\}_1^{\infty},
\end{align*}
并与交也采用类似记号。
\begin{definition}[集合序列的上极限与下极限]\label{definition:ch00-set-sequences}
在这种情形下，\textbf{上极限}与\textbf{下极限}的概念有时很有用：
\begin{align*}
\limsup E_n&=\bigcap_{k=1}^{\infty}\bigcup_{n=k}^{\infty}E_n,
&\liminf E_n&=\bigcup_{k=1}^{\infty}\bigcap_{n=k}^{\infty}E_n.
\end{align*}
\end{definition}
% SUPPLEMENT: id=v14-ch00-set-limits-remark; parent=ch00-set-sequences; kind=motivation+hypotheses; origin=independent-derivation
\begin{remark}
\textbf{动机与条件说明（重构新增）。}这两个集合描述成员资格随 $n$ 变化时的长期行为：上极限记录反复出现的点，下极限记录最终不再离开的点。指标为按通常次序排列的 $\mathbb{N}$，才有“从第 $k$ 项开始的尾部”和“最终”的含义；$n\geqslant k$ 的尾部不能误换成只检查有限个开头项。外层与内层运算的次序分别给出“对每个 $k$，存在 $n\geqslant k$”与“存在 $k$，对所有 $n\geqslant k$”，量词次序不可互换。无需假设集合列单调。改动有限个项不影响足够靠后的尾部，故不改变任一极限。删去有限个初始项并重新编号，虽会使尾部指标平移，但不会改变“属于无穷多个成员”或“只有有限多个例外”这两种性质，因此由下方证明的成员刻画，极限仍然不变。
\end{remark}
% END SUPPLEMENT: id=v14-ch00-set-limits-remark



\begin{remark}
读者可以验证
\begin{align*}
\limsup E_n&=\{x:x\in E_n\text{ 对无穷多个 }n\text{ 成立}\},\\
\liminf E_n&=\{x:x\in E_n\text{ 对除有限多个以外的所有 }n\text{ 成立}\}.
\end{align*}
\end{remark}
% SUPPLEMENT: id=v14-ch00-set-limits-proof; parent=ch00-set-sequences; kind=proof; origin=independent-derivation
\begin{proof}
\textbf{成员刻画的补写证明（重构新增）。}固定 $x$，令 $I_x=\{n\in\mathbb{N}:x\in E_n\}$。由并与交的定义，$x\in\limsup E_n$ 等价于：对每个 $k\in\mathbb{N}$，存在 $n\in I_x$ 且 $n\geqslant k$。若 $I_x$ 无限，它不可能包含于有限集 $\{1,\ldots,k-1\}$，所以这个条件成立。反之，若 $I_x$ 有限，$I_x=\varnothing$ 时取 $k=1$；否则取 $k=1+\max I_x$，均不存在这样的 $n$。这证明了上极限的刻画。

同样，$x\in\liminf E_n$ 等价于存在 $k$，使每个 $n\geqslant k$ 均属于 $I_x$。该条件蕴含 $\mathbb{N}\setminus I_x\subseteq\{1,\ldots,k-1\}$，因而例外指标有限。反之，若例外指标集为空，取 $k=1$；若它非空且有限，取 $k=1+\max(\mathbb{N}\setminus I_x)$。此时所有 $n\geqslant k$ 都满足 $x\in E_n$，得到下极限的刻画。
\end{proof}
% END SUPPLEMENT: id=v14-ch00-set-limits-proof


% SOURCE p=18; printed=3; object=ch00-set-differences; region=top

\begin{definition}[差与对称差]\label{definition:ch00-set-differences}
若 $E$ 和 $F$ 是集合，它们的\textbf{差}记为 $E\setminus F$：
\begin{align*}
E\setminus F&=\{x:x\in E\text{ 且 }x\notin F\},
\end{align*}
它们的\textbf{对称差}记为 $E\mathbin{\triangle}F$：
\begin{align*}
E\mathbin{\triangle}F&=(E\setminus F)\cup(F\setminus E).
\end{align*}
\end{definition}
% SUPPLEMENT: id=v14-ch00-difference-remark; parent=ch00-set-differences; kind=motivation+hypotheses; origin=independent-derivation
\begin{remark}
\textbf{动机与条件说明（重构新增）。}差集 $E\setminus F$ 表示从 $E$ 中删去同时属于 $F$ 的元素；必须保留 $x\in E$ 这个条件，否则会把与 $E$ 无关的元素也算入。这里不要求 $F\subseteq E$，因为 $F$ 在 $E$ 外的部分不影响删除结果。对称差记录两个集合成员资格不同的点：恰属于其中一个，而不同时属于两个。两部分 $E\setminus F$ 与 $F\setminus E$ 不交，因为同时属于二者会同时要求 $x\in E$ 与 $x\notin E$。差集有方向性，对称差则由交换两部分可见对 $E,F$ 对称。空集、相等集合均允许，例如 $E\mathbin{\triangle}E=\varnothing$。
\end{remark}
% END SUPPLEMENT: id=v14-ch00-difference-remark


% SOURCE p=18; printed=3; object=ch00-complement-demorgan; region=upper

\begin{definition}[补集]\label{definition:ch00-complement-demorgan}
当所讨论的所有集合均为某个固定集合 $X$ 的子集这一点已明确时，定义集合 $E$（在 $X$ 中）的\textbf{补集} $E^c$ 为
\begin{align*}
E^c&=X\setminus E.
\end{align*}
\end{definition}
% SUPPLEMENT: id=v14-ch00-complement-remark; parent=ch00-complement-demorgan; kind=motivation+hypotheses; origin=independent-derivation
\begin{remark}
\textbf{动机与条件说明（重构新增）。}补集是在固定讨论范围内表达“不属于 $E$”。指定同一个全集 $X$ 是必要的语义条件：同一 $E$ 在不同全集中的补集通常不同。$E\subseteq X$ 使补集成为 $X$ 中与 $E$ 正好分开的另一部分；每个 $x\in X$ 或属于 $E$，或不属于 $E$，故 $E\cup E^c=X$ 且 $E\cap E^c=\varnothing$。补集的上标记号是约定，不表示幂。以下 deMorgan 律两边的补集均相对于同一 $X$，不能在不同位置更换全集。
\end{remark}
% END SUPPLEMENT: id=v14-ch00-complement-remark



\begin{proposition}[deMorgan 律]\label{proposition:ch00-demorgan-laws}
此时有 \textbf{deMorgan 律}：
\begin{align*}
\left(\bigcup_{\alpha\in A}E_\alpha\right)^c&=\bigcap_{\alpha\in A}E_\alpha^c,
&\left(\bigcap_{\alpha\in A}E_\alpha\right)^c&=\bigcup_{\alpha\in A}E_\alpha^c.
\end{align*}
\end{proposition}
% SUPPLEMENT: id=v14-ch00-demorgan-proof; parent=ch00-complement-demorgan; kind=proof+motivation+hypotheses; origin=independent-derivation
\begin{proof}
\textbf{补写证明（重构新增）。}固定 $x\in X$。$x$ 属于并集的补集，当且仅当不存在 $\alpha\in A$ 使 $x\in E_\alpha$；这等价于对每个 $\alpha\in A$，$x\notin E_\alpha$，也就等价于对每个 $\alpha\in A$，$x\in E_\alpha^c$。由交的定义，这正是 $x\in\bigcap_{\alpha\in A}E_\alpha^c$，第一个等式成立。

$x$ 属于交集的补集，当且仅当“对所有 $\alpha\in A$，$x\in E_\alpha$”不成立；在通常的经典逻辑下，这等价于存在 $\alpha\in A$ 使 $x\notin E_\alpha$，也就是 $x\in\bigcup_{\alpha\in A}E_\alpha^c$。第二个等式成立。两边均为 $X$ 的子集，而任意 $x\in X$ 的成员资格已逐一对应，所以是集合相等。若 $A=\varnothing$，按固定全集中的空并、空交约定，第一个等式两边都等于 $X$，第二个等式两边都等于 $\varnothing$，仍然成立。
\end{proof}
\begin{remark}
\textbf{动机与条件说明（重构新增）。}deMorgan 律把“某项成立”的否定转为“每项均不成立”，把“每项成立”的否定转为“至少一项不成立”，是量词否定在集合运算中的形式。条件只有：$E_\alpha$ 均为固定全集 $X$ 的子集，各补集在 $X$ 中计算，并采用经典逻辑和空族运算约定。不需要有限指标集、可数性、互不相交或单调性；证明逐点处理任意指标族，未使用这些额外条件。
\end{remark}
% END SUPPLEMENT: id=v14-ch00-demorgan-proof


% SOURCE p=18; printed=3; object=ch00-products-relations; region=middle

\begin{definition}[Cartesian 积与关系]\label{definition:ch00-products-relations}
若 $X$ 和 $Y$ 是集合，它们的 \textbf{Cartesian 积} $X\times Y$ 是所有满足 $x\in X$、$y\in Y$ 的有序对 $(x,y)$ 构成的集合。从 $X$ 到 $Y$ 的\textbf{关系}是 $X\times Y$ 的一个子集。（若 $Y=X$，便称为 $X$ 上的关系。）若 $R$ 是从 $X$ 到 $Y$ 的关系，有时用 $xRy$ 表示 $(x,y)\in R$。
\end{definition}
% SUPPLEMENT: id=v14-ch00-relations-remark; parent=ch00-products-relations; kind=motivation+hypotheses; origin=independent-derivation
\begin{remark}
\textbf{动机与条件说明（重构新增）。}Cartesian 积把来自两个指定集合的选择放在两个不同位置中；使用有序对，才能区分输入和输出、关系的起点和终点。要求 $x\in X$、$y\in Y$ 分别限制两个位置的取值范围；任一因子为空时没有满足条件的有序对，所以积为空。把关系定义为积的任意子集，意味着它只筛选哪些配对成立，不要求每个 $x$ 有配对，也不要求配对唯一。$Y=X$ 只是在同一个集合内部比较元素的情形，不自动赋予关系反身性、对称性或其他性质。
\end{remark}
% END SUPPLEMENT: id=v14-ch00-relations-remark

最重要的几类关系如下：
\begin{itemize}
% SOURCE p=18; printed=3; object=ch00-equivalence-relations; region=middle
\item 等价关系。
\begin{definition}[等价关系与等价类]\label{definition:ch00-equivalence-relations}
$X$ 上的\textbf{等价关系}是 $X$ 上满足下列条件的关系 $R$：
\begin{align*}
&xRx\quad\text{对所有 }x\in X\text{ 成立},\\
&xRy\quad\text{当且仅当}\quad yRx,\\
&xRz\quad\text{只要存在某个 }y\text{ 使得 }xRy\text{ 且 }yRz.
\end{align*}
元素 $x$ 的\textbf{等价类}为 $\{y\in X:xRy\}$。
\end{definition}
$X$ 是这些等价类的不交并。
% SUPPLEMENT: id=v14-ch00-equivalence-partition; parent=ch00-equivalence-relations; kind=proof+motivation+hypotheses; origin=independent-derivation
\begin{proof}
\textbf{等价类分划的补写证明（重构新增）。}记 $[x]=\{y\in X:xRy\}$。由反身性，$x\in[x]$，所以每个元素都属于至少一个等价类，各类的并为 $X$。设 $[x]\cap[y]\ne\varnothing$，取 $z$ 属于这个交。于是 $xRz$ 且 $yRz$；由对称性 $zRy$，再由传递性得 $xRy$。若 $u\in[x]$，则 $xRu$，而 $yRx$ 由 $xRy$ 的对称性得到，再用传递性得 $yRu$，故 $u\in[y]$。这证明 $[x]\subseteq[y]$。交换 $x,y$，相同的两步推理给出 $[y]\subseteq[x]$，故两类相等。因此不同的等价类必不相交，与覆盖 $X$ 一起得到所述分划。这里“不交并”按不同的等价类取并，不能把同一个类按其中每个元素重复列出后仍称两两不交。
\end{proof}
\begin{remark}
\textbf{动机与条件说明（重构新增）。}等价关系用于把在某种标准下视为相同的元素归入同一类。三个条件各有作用：反身性使元素属于自己的类；对称性保证“相同”的比较不依赖方向；传递性保证通过中间元素建立的相同性能继续传递，正是上面证明相交类相等时使用的条件。这些条件不是任意关系自动具有的，也彼此不能替代：非空 $X$ 上的空关系对称且传递，但不反身；$\mathbb{R}$ 上的 $\leqslant$ 反身且传递，却不对称；$\{0,1,2\}$ 上由 $|x-y|\leqslant1$ 定义的关系反身且对称，但 $0R1$、$1R2$ 成立而 $0R2$ 不成立。定义中的所有变量都在同一个 $X$ 中，这使三个条件及等价类分划有共同的作用范围。
\end{remark}
% END SUPPLEMENT: id=v14-ch00-equivalence-partition

% SOURCE p=18; printed=3; object=ch00-ordering-forward-reference; region=middle-lower
\item 序关系。见\cref{section:ch00-orderings}。
% SOURCE p=18; printed=3; object=ch00-mappings-definition; region=lower
\item 映射。
\begin{definition}[映射]\label{definition:ch00-mappings-definition}
\textbf{映射} $f:X\to Y$ 是从 $X$ 到 $Y$ 的一个关系 $R$，满足对每个 $x\in X$，存在唯一的 $y\in Y$ 使得 $xRy$，此时写作 $y=f(x)$。映射有时也称为 map 或\textbf{函数}；我们通常将后一个名称留给 $Y$ 为 $\mathbb{C}$ 或其某个子集的情形。
\end{definition}
% SUPPLEMENT: id=v14-ch00-mapping-remark; parent=ch00-mappings-definition; kind=motivation+hypotheses; origin=independent-derivation
\begin{remark}
\textbf{动机与条件说明（重构新增）。}函数是在一般关系中挑出“每个输入确定一个输出”的特殊情形。“对每个 $x\in X$”明确规定整个定义域；“存在”排除某些输入没有输出的部分对应；“唯一”排除同一输入有多个输出。关系属于 $X\times Y$，则保证输出落在指定的 $Y$ 中。若去掉存在性，可以取非空 $X$ 上的空关系；若去掉唯一性，关系 $\{(0,0),(0,1)\}$ 就会允许同一输入有两个输出。这里不要求不同输入有不同输出，那是单射的额外条件。空定义域上到任意 $Y$ 都有唯一的空映射；若 $X$ 非空而 $Y$ 为空，则不存在映射。将“函数”一词通常留给数值输出，是本书的用语选择，并非关系定义本身的限制。
\end{remark}
% END SUPPLEMENT: id=v14-ch00-mapping-remark


\end{itemize}
% SOURCE p=18; printed=3; object=ch00-composition; region=lower

\begin{definition}[映射的复合]\label{definition:ch00-composition}
若 $f:X\to Y$ 和 $g:Y\to Z$ 是映射，用 $g\circ f$ 表示它们的\textbf{复合}：
\begin{align*}
g\circ f&:X\to Z,\qquad g\circ f(x)=g(f(x)).
\end{align*}
\end{definition}
% SUPPLEMENT: id=v14-ch00-composition-remark; parent=ch00-composition; kind=motivation+hypotheses; origin=independent-derivation
\begin{remark}
\textbf{动机与条件说明（重构新增）。}复合表示先执行 $f$ 再执行 $g$，所以写成 $g\circ f$，与操作的阅读顺序相反。对任意 $x\in X$，$f$ 的存在唯一性给出唯一 $f(x)\in Y$；$Y$ 是 $g$ 的定义域，因此又得到唯一 $g(f(x))\in Z$。这逐点证明了复合是 $X\to Z$ 的映射。定义中共用中间集合 $Y$ 是方便的类型匹配条件；更一般地，只需 $f(X)$ 包含于 $g$ 的定义域，就能逐点完成同样的构造。若这个包含条件失败，某些 $g(f(x))$ 就没有定义；但不需要单射、满射、连续性或其他结构。
\end{remark}
% END SUPPLEMENT: id=v14-ch00-composition-remark


% SOURCE p=18; printed=3; object=ch00-image-inverse-image; region=bottom

\begin{definition}[像与原像]\label{definition:ch00-image-inverse-image}
若 $D\subseteq X$、$E\subseteq Y$，定义 $D$ 在映射 $f:X\to Y$ 下的\textbf{像}与 $E$ 的\textbf{原像}为
\begin{align*}
f(D)&=\{f(x):x\in D\},\qquad f^{-1}(E)=\{x:f(x)\in E\}.
\end{align*}
\end{definition}
% SUPPLEMENT: id=v14-ch00-image-remark; parent=ch00-image-inverse-image; kind=motivation+hypotheses; origin=independent-derivation
\begin{remark}
\textbf{动机与条件说明（重构新增）。}像记录输入集合能产生哪些输出；原像记录哪些输入的输出满足给定条件。因此 $D\subseteq X$ 保证对其中元素能够计算 $f$，$E\subseteq Y$ 则把输出条件放在陪域内。原像定义中的 $x$ 在定义域 $X$ 中取值，准确说是 $\{x\in X:f(x)\in E\}$。此处 $f^{-1}$ 作用于子集，不要求 $f$ 可逆；它始终是从 $\mathscr{P}(Y)$ 到 $\mathscr{P}(X)$ 的映射。与之不同，逐点的逆映射 $Y\to X$ 需要后面定义的双射条件。空集也允许，且由“没有成员满足输入条件”直接得到 $f(\varnothing)=\varnothing$ 和 $f^{-1}(\varnothing)=\varnothing$。
\end{remark}
% END SUPPLEMENT: id=v14-ch00-image-remark


% SOURCE p=19; printed=4; object=ch00-inverse-image-operations; region=top

\begin{proposition}\label{proposition:ch00-inverse-image-operations}
容易验证，由第二个公式定义的映射 $f^{-1}:\mathscr{P}(Y)\to\mathscr{P}(X)$ 与并、交及补运算可交换：
\begin{align*}
f^{-1}\left(\bigcup_{\alpha\in A}E_\alpha\right)&=\bigcup_{\alpha\in A}f^{-1}(E_\alpha),\\
f^{-1}\left(\bigcap_{\alpha\in A}E_\alpha\right)&=\bigcap_{\alpha\in A}f^{-1}(E_\alpha),\\
f^{-1}(E^c)&=(f^{-1}(E))^c.
\end{align*}
\end{proposition}
% SUPPLEMENT: id=v14-ch00-inverse-image-proof; parent=ch00-inverse-image-operations; kind=proof+motivation+hypotheses; origin=independent-derivation
\begin{proof}
\textbf{补写证明（重构新增）。}固定 $x\in X$。它属于 $f^{-1}(\bigcup_{\alpha\in A}E_\alpha)$，当且仅当 $f(x)$ 属于某个 $E_\alpha$；这等价于 $x$ 属于某个 $f^{-1}(E_\alpha)$，即属于右端并集。它属于 $f^{-1}(\bigcap_{\alpha\in A}E_\alpha)$，当且仅当对每个 $\alpha\in A$，$f(x)\in E_\alpha$；这等价于对每个 $\alpha\in A$，$x\in f^{-1}(E_\alpha)$，即属于右端交集。最后，
\begin{align*}
x\in f^{-1}(E^c)
&\Longleftrightarrow f(x)\in Y\setminus E\\
&\Longleftrightarrow f(x)\notin E\\
&\Longleftrightarrow x\notin f^{-1}(E)
\Longleftrightarrow x\in X\setminus f^{-1}(E).
\end{align*}
第二个等价利用 $f(x)\in Y$，最后一个等价利用 $x\in X$。这证明了补集公式。任意 $x\in X$ 的成员资格均对应，故三式都是集合相等。对于空指标族，并式两边为空集；交式左边为 $f^{-1}(Y)=X$，右边为空族在 $X$ 中的交，也等于 $X$。
\end{proof}
\begin{remark}
\textbf{动机与条件说明（重构新增）。}原像只是把关于输出 $y$ 的成员判断代入 $y=f(x)$，因此保留“或”“且”“非”这三种逻辑组合。只需 $f$ 是在整个 $X$ 上有唯一取值的映射、$E_\alpha\subseteq Y$，并将左侧的补集相对于 $Y$、右侧的补集相对于 $X$ 计算。不要求单射或满射，也不限制指标族的基数。不同全集在这里各有明确角色；不能把两侧补集误在同一个未说明的范围内计算。
\end{remark}
% END SUPPLEMENT: id=v14-ch00-inverse-image-proof



\begin{remark}
（直接像映射 $f:\mathscr{P}(X)\to\mathscr{P}(Y)$ 与并运算可交换，但一般不与交或补运算可交换。）
\end{remark}
% SUPPLEMENT: id=v14-ch00-direct-image-proof; parent=ch00-inverse-image-operations; kind=proof; origin=independent-derivation
\begin{proof}
\textbf{直接像断言的补写证明与反例（重构新增）。}对任意集合族 $\{D_\alpha\}_{\alpha\in A}$，$y\in f(\bigcup_\alpha D_\alpha)$ 当且仅当存在 $x$ 及某个 $\alpha$，使 $x\in D_\alpha$ 且 $y=f(x)$；这等价于对某个 $\alpha$ 有 $y\in f(D_\alpha)$，即 $y\in\bigcup_\alpha f(D_\alpha)$。两边均为 $Y$ 的子集，故并运算可交换，空指标族时两边也均为空集。

对于交，取 $X=\{0,1\}$、$Y=\{a\}$，并令 $f(0)=f(1)=a$。若 $D_0=\{0\}$、$D_1=\{1\}$，则 $f(D_0\cap D_1)=f(\varnothing)=\varnothing$，而 $f(D_0)\cap f(D_1)=\{a\}$。对于补集，仍取此映射和 $D=\{0\}$，则 $f(X\setminus D)=\{a\}$，而 $Y\setminus f(D)=\varnothing$。两组不等式分别证明直接像一般不与交、补运算交换。
\end{proof}
% END SUPPLEMENT: id=v14-ch00-direct-image-proof



% SOURCE p=19; printed=4; object=ch00-mapping-types; region=upper-middle

\begin{definition}[定义域、值域及映射的类型]\label{definition:ch00-mapping-types}
若 $f:X\to Y$ 是映射，称 $X$ 为 $f$ 的\textbf{定义域}，$f(X)$ 为 $f$ 的\textbf{值域}。若 $f(x_1)=f(x_2)$ 仅在 $x_1=x_2$ 时成立，则称 $f$ 为\textbf{单射}；若 $f(X)=Y$，则称 $f$ 为\textbf{满射}；若 $f$ 既是单射又是满射，则称它为\textbf{双射}。若 $f$ 是双射，它有\textbf{逆映射} $f^{-1}:Y\to X$，使得 $f^{-1}\circ f$ 和 $f\circ f^{-1}$ 分别为 $X$ 和 $Y$ 上的恒等映射。
\end{definition}
% SUPPLEMENT: id=v14-ch00-mapping-types-remark; parent=ch00-mapping-types; kind=proof+motivation+hypotheses; origin=independent-derivation
\begin{remark}
\textbf{动机与条件说明（重构新增）。}定义域是允许的全部输入；陪域 $Y$ 是预先指定的输出范围，而值域 $f(X)$ 是实际达到的部分。单射表示不同输入不会合并，满射表示陪域没有遗漏，双射同时满足两者，因此可完整恢复输入。满射性依赖所指定的陪域：同一取值规则若把陪域改为实际值域，就成为满射。单射可等价写成 $x_1\ne x_2\Rightarrow f(x_1)\ne f(x_2)$，这是对应蕴含的逆否命题。双射存在逆映射时，满射性负责对每个 $y$ 找到原像，单射性负责原像唯一，两个条件分工不同。
\end{remark}
\begin{proof}
\textbf{逆映射存在唯一性的补写证明（重构新增）。}设 $f$ 为双射。对每个 $y\in Y$，满射性给出 $x\in X$ 使 $f(x)=y$；若 $f(x')=y$，单射性给出 $x'=x$。故可以将这个唯一的 $x$ 定义为 $g(y)$，得到映射 $g:Y\to X$。对 $x\in X$，$x$ 正是 $f(x)$ 的唯一原像，所以 $g(f(x))=x$；对 $y\in Y$，构造本身给出 $f(g(y))=y$，这就是两条恒等映射关系。若 $h:Y\to X$ 也满足这两条关系，则对任意 $y\in Y$，$f(h(y))=y=f(g(y))$；单射性给出 $h(y)=g(y)$，故逆映射唯一。空集情形也包含在这些逐点陈述中：双射的一个集合为空时另一个也必须为空，逆映射就是唯一的空映射。
\end{proof}
% END SUPPLEMENT: id=v14-ch00-mapping-types-remark


\begin{definition}[映射的限制]\label{definition:ch00-mapping-restriction}
若 $A\subseteq X$，用 $f|A$ 表示 $f$ 在 $A$ 上的\textbf{限制}：
\begin{align*}
(f|A)&:A\to Y,\qquad (f|A)(x)=f(x)\quad\text{对 }x\in A.
\end{align*}
\end{definition}
% SUPPLEMENT: id=v14-ch00-restriction-remark; parent=ch00-mapping-types; kind=motivation+hypotheses; origin=independent-derivation
\begin{remark}
\textbf{动机与条件说明（重构新增）。}限制保留原有取值规则，只缩小允许的输入范围。$A\subseteq X$ 确保每个新的输入原本就在 $f$ 的定义域中，所以其输出仍存在、唯一且属于 $Y$，这也验证了 $f|A$ 是映射。$A=X$ 时得到原映射，$A=\varnothing$ 时得到空映射，均无需排除。限制通常改变实际值域，而定义中的陪域仍为 $Y$，所以不能自动把限制称为满射。
\end{remark}
% END SUPPLEMENT: id=v14-ch00-restriction-remark


% SOURCE p=19; printed=4; object=ch00-sequences-subsequences; region=middle

\begin{definition}[序列与子列]\label{definition:ch00-sequences-subsequences}
集合 $X$ 中的\textbf{序列}是从 $\mathbb{N}$ 到 $X$ 的映射。（我们也用\textbf{有限序列}表示从 $\{1,\ldots,n\}$ 到 $X$ 的映射，其中 $n\in\mathbb{N}$。）若 $f:\mathbb{N}\to X$ 是序列，而 $g:\mathbb{N}\to\mathbb{N}$ 满足只要 $n<m$ 就有 $g(n)<g(m)$，则称复合映射 $f\circ g$ 为 $f$ 的\textbf{子列}。
\end{definition}
% SUPPLEMENT: id=v14-ch00-subsequence-remark; parent=ch00-sequences-subsequences; kind=motivation+hypotheses; origin=independent-derivation
\begin{remark}
\textbf{动机与条件说明（重构新增）。}把序列视为映射，可以保留次序和重复次数；只看值域会丢失这些信息。$\mathbb{N}$ 在本书从 $1$ 起，有限序列的 $n\in\mathbb{N}$ 表示这里约定其长度为正，并非断言空序列在其他约定中不能定义。子列的 $g:\mathbb{N}\to\mathbb{N}$ 保证抽取无限多个有合法指标的项；严格递增保证保持原顺序且不重复使用同一指标。由于 $g(1)\geqslant1$，且整数严格递增给出 $g(n+1)\geqslant g(n)+1$，归纳得 $g(n)\geqslant n$，所以选取指标趋于无穷。若仅要求非递减，则 $g(n)=1$ 会不断重复第一项，不符合这里的子列定义。不同指标处的函数值仍可相同，严格性约束的是指标而非数值。
\end{remark}
% END SUPPLEMENT: id=v14-ch00-subsequence-remark


\begin{remark}
通常不严格区分序列及其值域，这往往也很方便；其值域是 $X$ 的以 $\mathbb{N}$ 为指标集的子集。因此，若 $f(n)=x_n$，便称序列为 $\{x_n\}_1^\infty$；究竟指从 $\mathbb{N}$ 到 $X$ 的映射，还是指 $X$ 的一个子集，由上下文即可明确。
\end{remark}


% SOURCE p=19; printed=4; object=ch00-arbitrary-cartesian-product; region=middle-lower
前面定义了两个集合的 Cartesian 积。同样，可以用有序 $n$ 元组定义 $n$ 个集合的 Cartesian 积。然而，对无限集合族而言，这个定义很不方便，因此采用以下方式。
\begin{definition}[集合族的 Cartesian 积]\label{definition:ch00-arbitrary-cartesian-product}
若 $\{X_\alpha\}_{\alpha\in A}$ 是带指标的集合族，它们的 \textbf{Cartesian 积} $\prod\limits_{\alpha\in A}X_\alpha$ 是所有满足 $f(\alpha)\in X_\alpha$（对每个 $\alpha\in A$）的映射 $f:A\to\bigcup_{\alpha\in A}X_\alpha$ 构成的集合。
\end{definition}
% SUPPLEMENT: id=v14-ch00-product-remark; parent=ch00-arbitrary-cartesian-product; kind=motivation+hypotheses; origin=independent-derivation
\begin{remark}
\textbf{动机与条件说明（重构新增）。}用函数记录“在每个位置选了什么”，就能统一有限元组与任意指标的元组。$A$ 是指标集合，各 $X_\alpha$ 是该位置允许的取值集合；取函数的陪域为所有因子的并，使每个可能值都有容身之处，而逐项条件 $f(\alpha)\in X_\alpha$ 进一步排除选错位置的值。只要求 $f(\alpha)$ 属于总并是不够的。若某个因子为空而指标确实存在，就没有合法选择，乘积为空；若 $A=\varnothing$，唯一的空函数满足全部条件，乘积为只含空函数的单元素集合。对任意无限指标族，“各因子非空便能同时选出一个元素”涉及选择公理；乘积的定义本身并未自动证明其非空。
\end{remark}
% END SUPPLEMENT: id=v14-ch00-product-remark


\begin{remark}
（应当注意，然后立即忘掉：当 $A=\{1,2\}$ 时，先前对 $X_1\times X_2$ 的定义在集合论意义上不同于现在对 $\prod\limits_1^2X_j$ 的定义。事实上，后一个概念依赖映射，而映射又是用前一个概念定义的。）
\end{remark}

\begin{definition}[投影与坐标]\label{definition:ch00-projection-coordinate}
若 $X=\prod\limits_{\alpha\in A}X_\alpha$ 且 $\alpha\in A$，定义第 $\alpha$ 个\textbf{投影}或\textbf{坐标映射} $\pi_\alpha:X\to X_\alpha$ 为 $\pi_\alpha(f)=f(\alpha)$。我们也常常分别用 $x$、$x_\alpha$ 代替 $f$、$f(\alpha)$，并称 $x_\alpha$ 为 $x$ 的第 $\alpha$ 个\textbf{坐标}。
\end{definition}
% SUPPLEMENT: id=v14-ch00-projection-remark; parent=ch00-projection-coordinate; kind=motivation+hypotheses; origin=independent-derivation
\begin{remark}
\textbf{动机与条件说明（重构新增）。}投影从一个完整的选择中读取指定位置的数据。必须有 $\alpha\in A$，否则函数 $f$ 没有这一指标的值；必须让输入 $f$ 属于乘积，才能由乘积条件保证 $f(\alpha)\in X_\alpha$。函数在一个指标上的值唯一，所以 $\pi_\alpha$ 是良定义的映射。不需要所有因子相同，也不需要额外的拓扑或代数结构。乘积为空时，投影是空定义域上的映射；定义没有自动保证投影满射。比如 $A=\{1,2\}$、$X_1=\{0\}$、$X_2=\varnothing$，则 $\pi_1$ 的定义域为空，其值域不等于 $X_1$。
\end{remark}
% END SUPPLEMENT: id=v14-ch00-projection-remark



% SOURCE p=19; printed=4; object=ch00-function-set-power; region=lower

\begin{definition}[映射集合的幂记号]\label{definition:ch00-function-set-power}
若所有 $X_\alpha$ 均等于某个固定集合 $Y$，便用 $Y^A$ 表示 $\prod\limits_{\alpha\in A}X_\alpha$：
\begin{align*}
Y^A&=\text{从 }A\text{ 到 }Y\text{ 的所有映射构成的集合}.
\end{align*}
若 $A=\{1,\ldots,n\}$，则将 $Y^A$ 记为 $Y^n$，并可将其与由 $Y$ 中元素组成的有序 $n$ 元组的集合等同起来。
\end{definition}
% SUPPLEMENT: id=v14-ch00-function-power-remark; parent=ch00-function-set-power; kind=motivation+hypotheses; origin=independent-derivation
\begin{remark}
\textbf{动机与条件说明（重构新增）。}幂记号 $Y^A$ 强调每个指标都使用同一个值集合 $Y$；若各因子不同，就不能直接用同一 $Y$ 代替而丢掉逐项限制。对于 $A=\{1,\ldots,n\}$，函数 $f$ 对应有序元组 $(f(1),\ldots,f(n))$；反过来，元组 $(y_1,\ldots,y_n)$ 确定唯一函数 $j\mapsto y_j$。两个构造互相抵消，因此这里的“等同”是自然的双射识别，并非断言采用不同集合论编码的对象逐字相同。任意 $A$ 下仍须保留指标信息，而不是只保留函数值组成的无序集合。
\end{remark}
% END SUPPLEMENT: id=v14-ch00-function-power-remark


\end{document}
